% !TeX TS-program = xelatex

\documentclass{resume}
\ResumeName{练嘉铖}

% 如果想插入照片，请使用以下两个库。
% \usepackage{graphicx}
% \usepackage{tikz}

\begin{document}

\ResumeContacts{
  18962562186,%
  \ResumeUrl{mailto:23009200100@stu.xidian.edu.cn}{23009200100@stu.xidian.edu.cn},%
  \ResumeUrl{https://github.com/WestNorth-Wolf}{https://github.com/WestNorth-Wolf},%
  江苏·张家港%
}

% 如果想插入照片，请取消此代码的注释。
% 但是默认不推荐插入照片，因为这不是简历的重点。
% 如果默认的照片插入格式不能满足你的需求，你可以尝试调整照片的大小，或者使用其他的插入照片的方法。
% 不然，也可以先渲染 PDF 简历，然后用其他工具在 PDF 上叠加照片。
% \begin{tikzpicture}[remember picture, overlay]
%   \node [anchor=north east, inner sep=1cm]  at (current page.north east) 
%      {\includegraphics[width=2cm]{image.png}};
% \end{tikzpicture}

\ResumeTitle


\section{教育经历}
\ResumeItem
[XiDian University 本科三年级]
{西安电子科技大学}
[\textnormal{计算机科学与技术|} 本科三年级在读]
[2023.09—2027.06(预计)]

\textbf{均分：94.7，排名：1/410，CET4：599，CET6：573。}(截至2025.8)

连续两年获得国家奖学金，优秀学生标兵称号，2025年获得比亚迪奖学金。



\section[技术能力]{技术能力}
\begin{itemize}
  \item \textbf{编程语言}: 对编程语言有较强的适应能力。熟悉C/C++(CMake)、Python。了解Java、Matlab。
  \item \textbf{工作流}: Linux, VSCode, JetBrains IDE, Git, GitHub。
  \item \textbf{嵌入式}: 熟悉嵌入式开发，主要使用STM32F4系列单片机基于FreeRTOS操作系统进行开发，实现传感器数据采集、控制逻辑、电机控制、与上位机通信等任务。了解嵌入式Linux操作系统、ROS2机器人操作系统。
  \item \textbf{其他}: 熟悉项目管理基础知识与常用工具，开发并实践简单项目管理框架。 
\end{itemize}

\section{竞赛获奖}

\ResumeItem
{二〇二五年全国大学生电子设计竞赛}
[全国二等奖]
[2025.08]

\ResumeItem
{第二十四届全国大学生机器人大赛RoboMaster2025机甲大师超级对抗赛·全国赛}
[三等奖]
[2025.08]

\ResumeItem
{第二十三届全国大学生机器人大赛RoboMaster2024机甲大师超级对抗赛·全国赛}
[二等奖]
[2024.08]

\ResumeItem
{第十七届全国大学生数学竞赛(非数学A类)}
[一等奖]
[2025.12]

\ResumeItem
{第三十九届全国中学生物理竞赛(江苏赛区)}
[一等奖]
[2022.09]

\section{项目经历}

\ResumeItem{IRobot战队成员}
[嵌入式组组员(2024/2025赛季)、项目管理(2026赛季)]
[2023.09—至今]
\begin{itemize}
  \item 嵌入式硬件：掌握嘉立创EDA的使用，能够设计简单电路板，负责英雄机器人的硬件电路部分。
  \item 嵌入式软件：使用嵌入式C/C++编程，进行STM32F4系列单片机的开发，实现机器人的电控控制及通信。
  \item 英雄组负责人/操作手：对机器人的研发进度进行跟踪管理，负责组织进行调试与测试，对机器人上场表现负责。
  \item 项目管理：基于OKR方法，制定项目计划，分配任务，跟踪进度，协调合作，发现风险并处理，确保研发进度正常。
\end{itemize}

\ResumeItem{二〇二五年全国大学生电子设计竞赛}
[嵌入式、队长]
[2025.08]
\begin{itemize}
  \item 嵌入式开发：负责C题（基于单目视觉的目标物测量装置）中嵌入式系统的软硬件开发。
  \item 综合测评：负责电赛综合测评环节的主要电路设计、焊接与调试工作。
  \item 队长：负责整个队伍的项目管理，协调任务分配和团队合作，参与视觉算法研发与调试。
\end{itemize}

\ResumeItem{智能巡逻安防机器人}
[大学生创新创业项目]
[2023.12-2025.06]
\begin{itemize}
  \item 负责机器人的嵌入式开发部分。
  \item 参与视觉、导航部分算法开发。学习ROS2机器人操作系统，了解SLAM与导航算法。
  \item 使用深度学习模型(Vision Transformer)进行人物检测与识别。
\end{itemize}

% \ResumeItem{人脸识别项目}
% [个人学习项目]
% \begin{itemize}
%   \item 使用 MediaPipe 和 OpenCV 完成人脸检测，并基于关键点进行几何对齐与标准化裁剪。
%   \item 采用 Vision Transformer (ViT) 作为骨干网络提取面部语义特征，充分利用其全局感受野捕捉面部细节关联，提升识别准确率。
%   \item 基于 TensorFlow 进行微调训练，在测试集上取得了高精度的识别效果。
% \end{itemize}

\section{个人陈述}

\begin{itemize}
  \item 本人成绩优异、学习能力出色、自驱力强，具有优秀的沟通和团队协作能力。
  \item 本人积极参与各类竞赛与工程项目，积累了丰富的实践经验，具备解决复杂工程问题的能力。
  \item 本人对机器人和人工智能技术充满热情，自初中起就关注相关领域的发展，看好具身智能的广阔前景，渴望在这一领域做出贡献。
  \item 本人符合参与本次活动的基本条件。不足之处在于对深度学习框架的学习和使用比较有限，近期也正在系统学习相关知识，希望能通过这次冬令营提升自己。
  \item 南京大学智能科学与技术学院（苏州）离我的家乡较近。本次活动对我来说是一个难得的学习和成长机会，希望能够参加此次科学营。
\end{itemize}


\end{document}
